\documentclass[12pt,a4paper]{article}
\usepackage[utf8]{inputenc}
\usepackage{ctex}
\usepackage{amsmath}
\usepackage{amssymb}
\usepackage{geometry}
\usepackage{booktabs}

\geometry{left=2.5cm,right=2.5cm,top=3cm,bottom=3cm}

\title{\textbf{Optimal Resizable Arrays 论文评述}}
\author{张艺博 \and 徐黄浩}
\date{\today}

\begin{document}

\maketitle

\begin{abstract}
本文评述 Robert E. Tarjan 与 Uri Zwick 的论文 \emph{Optimal resizable arrays}。该论文研究如何在保持数组随机访问效率的同时，降低动态扩缩容带来的额外空间开销。本文将围绕问题模型、相关工作、核心数据结构、复杂度证明以及方法局限展开分析。
\end{abstract}

\newpage
\tableofcontents
\newpage

\section{引言与论文定位}

数组和可变数组是广泛应用的数据结构，可变数组的内存效率一直是研究者的研究热点。本文所研究的可变数组指的是，可以在末尾添加或者删除元素的数组。为此，很自然想到，如何解决可变长数组的空间分配问题？一个典型场景是：内存已分配的空间已经被填满，此时需要插入新元素，却不能直接在原有内存块末尾追加，因为紧邻的内存可能已被其他对象占用。因此，必须重新申请一块更大的内存，将原有元素全部复制过去，再释放旧块。

这一问题可以从\textbf{机制}与\textbf{策略}两个层面加以分离与理解。

\paragraph{机制层面}
\begin{enumerate}
    \item 系统提供固定大小的存储块（block）；
    \item 存储块可以存放实际数据，也可以存放指向其他块的指针；
    \item Grow 操作：申请新块，必要时进行数据迁移；
    \item Shrink 操作：删除末尾元素，释放不再需要的空块。
\end{enumerate}

\paragraph{策略层面}
\begin{enumerate}
    \item 块的大小如何选取？
    \item 使用多少层指针？Grow 过程中块的大小如何变化？
    \item Shrink 的触发条件是什么？何时进行全局重建（rebuild）？
\end{enumerate}

将机制与策略解耦，有助于清晰地分析和设计可变长数组的数据结构。

\paragraph{经典方法：翻倍策略（Doubling）}
工业界广泛采用的成熟做法是几何增长策略，其核心规则为：当数组满时，分配一块大小为当前数组两倍的连续内存，并将所有元素复制过去。这一策略的摊还代价为 $O(1)$，
随机访问同样为 $O(1)$ 最坏情况时间。标准倍增算法可以达到 $O(N)$ (更确切说，应该是 $ N+O(N) $)空间开销和$O(1)$ 的时间开销，但其最坏时刻的空间浪费可达 $50\%$。进一步，我们可以将这些策略总结成如下规律(再后面第三节传统文献调研部分有详细说明)：
假设每次数组满后，增长因子为 $ (1+\alpha) $，那么新数组刚扩展完毕时浪费率为 $ \frac{\alpha}{1+\alpha} $，但是， $ \alpha $ 也不能无限小，当其越小时，时间开销越大，因为平均复制次数增多了，假设数组容量为$ C $后触发 grow，新容量变为$ C' = (1+\alpha)C $，每个操作的成本为 $ 1/\alpha $

\paragraph{本文的目标}
能否将空间占用压缩至 $N+o(N)$，同时仍然保持 $O(1)$ 的随机访问时间？


\section{问题定义与理论模型}


\section{文献调研与相关工作}

在进行相关工作阐述之前，我们先解释一下摊还分析(Amortized Analysis)。

摊还分析是一种分析操作序列平均代价的技术，由于难以衡量单次操作的最坏代价，所以我们从总代价的角度出发，使用总代价除以操作次数，得到平均每次操作花费代价。

具体而言，有三种分析方法：

1. 聚合分析：直接计算总代价，除以操作次数
2. 记账法
3. 势能法

本文中使用的主要是第一种方法。

\subsection{传统工作}
首先介绍论文中提到的相关工作

\paragraph{朴素方法}
很自然想到，最直接实现可变长数组的方案就是维护一个大小恰好等于当前元素数 $ N $ 的固定数组，每次 grow 或者 shrink 时都触发全量重建

它的存储空间显然是 $N+O(1)$，可以达到理论最优，但是 grow 或者 shrink 的摊还代价为 $\Theta(N)$，并且临时空间最大为 $2N+O(1)$

\paragraph{Geometric Expansion} 第二种传统方法就是前文提到过的\textbf{几何增长策略}，用空间换取时间，通过预分配额外容量降低重建成本。它的具体方法如下
假设增长因子为$1+\alpha$，初始分配大小为$C$的数组，当$N$个元素填满后，分配大小为$(1+\alpha)N$的新数组，复制全部元素，当元素数降至 $\frac{N}{1+\alpha}$ 以下时，收缩至 $\frac{N}{1+\alpha}$

使用摊还分析该数据结构，它的存储空间为$O(N)$，grow摊还代价为$N/(\alpha N) = 1/\alpha$

可以看到，如果我们增大$\alpha$，虽然摊还代价小了，但是我们总的存储空间增加，因此，传统方法的局限就在于此，难以平衡空间和时间。根本原因在于粒度太粗，也就是现在的空间划分只是计算的平均空间，后续作者针对这种不足提出了本文的改进。

\subsection{其他传统工作}
这里介绍本文中未提及的其他工作

实际上，从工程角度出发，很多语言都有自己的方法，但是绝大部分都是遵循的前面提到的第二种方法：几何增长策略，只不过增长因子不同。例如：C++ std::vector使用$\alpha=2$，Java Arraylist使用$\alpha=1.5$，python list使用1.125。
这里有一个非常有意思的点，就是Java中选择增长因子为1.5。我们考虑一个关键的工程因素，也就是释放后的旧块未来能否被新块利用？如果说旧块累积总容量超过了需要的新块，那么内存分配器就可以分配，否则，内存很容易满（更容易理解的说法是
如果我下一次只需要临时用一点内存，但是分配的新块如果很大，结果没有足够的内存，但实际上我们只需要分配一点内存就能满足，这说明$\alpha$不能过大。

此外，链表与分块链表也是解决动态数组的一种策略，考虑一般的分块链表，每个节点为一个大小为 $B$ 的数组块，虽然其能够做到空间$N + O(N/B)$，但是因为它无法做到随机访问（访问时间为$O(N)$），不能成为动态数组的替代。

Phil等人提出的VList对上面的分块链表做了改进，将每个节点的Block大小设置为可变的，大小按照几何比例变化（1,2,4,8...），这其实是一个很好的平衡，它使得访问时间变为$O(log N)$，但是最坏存储空间浪费为$2N+O(log N)$（这是因为刚扩容的时候申请的空闲块大小为前面的两倍），这是其不足的一点，此外，

Bagwell等人提出的VArray继续延伸VList思想，

\section{本文核心方法}


\section{正确性证明与理论分析}


\section{技术直觉与方法评价}


\section{总结}


\end{document}
